% ============================================================
%  Chapter 8 — Medical Robots
%  Robot World: Meet the Machines That Help Us
% ============================================================

\chapter{Medical Robots}

\chapteropening{%
  Robots that save lives. From operating theatres to hospital corridors,
  robots are helping doctors, nurses, and patients every day.
}
\chapterbadge{Sparks: Precision That Saves Lives}

% --- Introduction ---
\section*{Robots in the Hospital}

Hospitals need extreme precision and cleanliness. Robots can
hold instruments perfectly still, deliver medicines without getting
tired, and even help people learn to walk again.

\character{Sparks}{I work inside the human body! I am so small and
  precise that I can stitch a wound only a few millimetres wide.
  My steady hands never tremble.}

\sectionrule

% --- Surgical robots ---
\section*{Surgical Robots}

The most famous is the \textbf{\term{da Vinci Surgical System}}. A surgeon
sits at a console and controls tiny robotic arms that operate through
incisions smaller than a centimetre.

Advantages:
\begin{itemize}
  \item Smaller cuts mean less pain and faster healing.
  \item The camera gives a 3D magnified view inside the body.
  \item Robot arms filter out any hand tremors.
  \item Surgeons can operate remotely (\term{telesurgery}).
\end{itemize}

\begin{funfact}
  The da Vinci system has been used in over 12 million surgeries
  worldwide since it was approved in 2000.
\end{funfact}

\sectionrule

% --- Pharmacy and delivery robots ---
\section*{Hospital Helper Robots}

\begin{itemize}
  \item \textbf{\term{Pharmacy robots}} count and package pills automatically,
    reducing mistakes.
  \item \textbf{\term{Delivery robots}} carry blood samples, meals, and laundry
    through hospital corridors using maps and sensors.
  \item \textbf{\term{Disinfection robots}} blast rooms with ultraviolet light
    to kill germs after a patient leaves.
\end{itemize}

\begin{picturethis}
  Imagine a tiny robot postman rolling down the hospital corridor at night,
  delivering medicine to the right ward without ever getting a room number
  wrong. That is a hospital delivery robot.
\end{picturethis}

\sectionrule

% --- Rehabilitation robots ---
\section*{Rehabilitation Robots}

When someone has a stroke or a spinal injury, they may need to re-learn
how to walk or move their arms. \term{Exoskeletons} are wearable robots
that support the body and guide movements during therapy.

\begin{bigidea}
  Rehabilitation robots do not replace physiotherapists \textemdash{}
  they help patients practise for longer and track their progress
  more precisely.
\end{bigidea}

\sectionrule

% --- Micro-bots ---
\section*{Micro-Bots and Nano-Bots}

Scientists are developing robots so small they could travel through
blood vessels to deliver drugs directly to a tumour or unblock an artery.
These are still mostly experimental, but the future is exciting.

\character{Sparks}{One day, robots like me could swim through your
  bloodstream and fix problems before you even feel ill!}

\sectionrule


% --- How surgical robot works ---
\begin{quickmap}
  \textbf{How a robotic surgery works:}
  \begin{enumerate}
    \item The surgeon sits at a console with 3D screens and hand controllers.
    \item Tiny robotic arms are inserted through small incisions in the patient.
    \item A camera arm shows a magnified, high-definition view inside the body.
    \item The surgeon moves the controllers; the robot copies the movements exactly.
    \item Software filters out any hand tremor, keeping instruments perfectly steady.
    \item When finished, the tiny cuts heal much faster than a large open wound.
  \end{enumerate}
\end{quickmap}

\sectionrule

% --- Diagram ---
\begin{diagram}[Surgical robot setup: surgeon at console, robot arms at patient.]
  \begin{tikzpicture}[>=Stealth]
    % Surgeon console
    \draw[fill=SteelBlue!15,draw=SteelBlue,line width=1pt,rounded corners=4pt]
      (0,0) rectangle (2.5,2.5);
    \node[font=\small\bfseries,align=center] at (1.25,2) {Surgeon};
    \node[font=\small,align=center] at (1.25,1.2) {Console};
    \node[font=\scriptsize,align=center] at (1.25,0.5) {(3D view +\\controls)};
    % Connection line
    \draw[RoboGreen!70!black,line width=2pt,->] (2.5,1.25) -- (5,1.25);
    \node[font=\scriptsize\itshape,above] at (3.75,1.35) {commands};
    % Robot arms
    \draw[fill=RoboGreen!15,draw=RoboGreen!80!black,line width=1pt,rounded corners=4pt]
      (5,0) rectangle (8,2.5);
    \node[font=\small\bfseries,align=center] at (6.5,2) {Robot Arms};
    \node[font=\small,align=center] at (6.5,1.2) {+ Camera};
    \node[font=\scriptsize,align=center] at (6.5,0.5) {(at patient)};
    % Feedback arrow
    \draw[CoralRed!70,line width=1.5pt,->] (5,0.5) -- (2.5,0.5);
    \node[font=\scriptsize\itshape,below] at (3.75,0.35) {video feed};
  \end{tikzpicture}
\end{diagram}

\sectionrule

% --- Try It ---
\begin{tryit}
  Test your own hand steadiness:
  \begin{enumerate}
    \item Draw a straight line on paper (use a ruler).
    \item Now try to trace over it freehand with a pen.
    \item Try again while someone gently shakes the table.
  \end{enumerate}
  Surgical robots filter out shaking like this \textemdash{}
  the surgeon can move, but the instrument stays perfectly steady.
\end{tryit}

\sectionrule


\begin{quickcheck}
  \begin{enumerate}
    \item How does the da Vinci system help surgeons?
    \item What is an exoskeleton and who does it help?
    \item Why might a micro-bot be useful inside the human body?
  \end{enumerate}
\end{quickcheck}

\sectionrule
% --- New Words ---
\begin{newwords}
  \begin{description}[labelwidth=3.8cm, leftmargin=4.0cm]
    \item[Surgical robot] A robot system that helps surgeons perform operations.
    \item[Telesurgery] Surgery performed remotely using a robot.
    \item[Exoskeleton] A wearable robotic frame that supports or strengthens the body.
    \item[Micro-bot] A very small robot, often smaller than a fingernail.
    \item[Nano-bot] A theoretical robot at the scale of molecules.
    \item[Ultraviolet (UV)] Invisible light used to kill bacteria and viruses.
  \end{description}
\end{newwords}
